\section{项目概述与设计目标}

\subsection{赛题要求对齐}

结合赛题题面要求，本项目的说明书需要覆盖以下四类正式内容：
\begin{itemize}
    \item CPU 架构设计说明书；
    \item Verilog/VHDL 建模规范；
    \item 验证计划与测试报告；
    \item FPGA 适配指南。
\end{itemize}

围绕上述要求，本设计采用五级顺序流水结构，以 \texttt{RV32I + Zicsr} 为核心冻结子集，
围绕最小 SoC 和 FPGA 原型实现建立完整工程链路。说明书编排上不再采用泛化的工程周报式结构，
而是按赛题要求把“架构说明、建模规范、验证报告、FPGA 适配”作为正文主线，便于评审直接对照核查。
同时，本文采用更严格的对齐口径：赛题中出现的能力、流程、材料和评分点均视为赛方希望看到的目标项，
因此正文里的每一个“亮点”都必须能够回指到具体赛题要求，而不是脱离题面单独叙述。

\subsection{设计目标}

本项目在初赛阶段设置了五项核心目标：
\begin{enumerate}
    \item \textbf{完整 CPU 基线}：围绕 \texttt{RV32I/RV64I} 题面要求构建功能完整的处理器基线，并明确冻结提交口径与扩展验证口径；
    \item \textbf{五级流水与冒险闭环}：完成 IF/ID/EX/MEM/WB 五级流水组织，并对数据冒险、控制冒险和结构冒险给出工程化实现路径；
    \item \textbf{全流程工程链}：建立从 HDL 建模、仿真验证、软件运行到 FPGA 原型适配的统一链路；
    \item \textbf{性能优化可说明}：把已落地的前递、最小停顿与控制流优化边界写清楚，使性能亮点直接对应题面要求；
    \item \textbf{交付与后续任务清晰}：把简单应用、演示材料、源码打包、板级闭环等未完成项显式列入后续任务，而不是在说明书中隐去。
\end{enumerate}

\begin{table}[H]
    \centering
    \small
    \caption{赛题目标、本文亮点与当前口径}
    \begin{tabularx}{0.95\textwidth}{L{0.26\textwidth} Y L{0.14\textwidth}}
        \toprule
        赛题目标 & 本文亮点与当前口径 & 当前状态 \\
        \midrule
        基于 \texttt{RV32I/RV64I} 的完整 CPU & 以 \texttt{RV32I + Zicsr} 作为冻结提交口径，并用 \texttt{RV32/RV64} 双 XLEN baseline / full-ui 结果说明工程具备更广验证能力 & 已覆盖 \\
        五级流水与三类冒险处理 & 以五级顺序流水、两级前递、load-use 最小停顿、redirect/trap 统一控制流为主线组织正文 & 已覆盖 \\
        HDL 建模到 FPGA 的全流程 & 说明书串联 RTL、TestBench、软件镜像、CoreMark、Vivado \texttt{impl50} 与 FPGA-like probe 证据链 & 已覆盖 \\
        至少两项性能优化技术 & 当前明确写出数据旁路与精确停顿控制；控制流前瞻机制作为继续对齐赛题的后续增强方向单列说明 & 部分覆盖 \\
        测试程序、简单应用与功能演示 & 当前已形成 smoke 固件、\texttt{riscv-tests} 与 CoreMark 路径；排序/矩阵类简单应用与演示视频仍需补齐 & 待补齐 \\
        完整交付件与评分项对齐 & 正文按赛题结构重排，并在后续工作章节中单列源码包、注释率、Word/PDF、PPT、视频等交付缺口 & 待补齐 \\
        \bottomrule
    \end{tabularx}
\end{table}

\begin{table}[H]
    \centering
    \small
    \caption{工程特征与当前边界}
    \begin{tabularx}{0.95\textwidth}{L{0.22\textwidth} Y}
        \toprule
        主题 & 当前口径 \\
        \midrule
        冻结主口径 & 采用 \texttt{RV32I + Zicsr} 五级顺序流水，作为当前对外性能与提交主叙述路径 \\
        扩展验证能力 & 当前工程验证已覆盖 \texttt{RV32/RV64} 双 XLEN，并具备 baseline 与更广 \texttt{full-ui} 覆盖 \\
        SoC 组织 & 采用同步指令存储、数据存储与最小 MMIO 外设构成闭环，优先保证仿真、回归与 FPGA 路径一致性 \\
        机器态支持 & 当前重点覆盖 CSR、trap、\texttt{mret} 和 timer interrupt 等比赛阶段确有需求的机器态能力 \\
        原型实现目标 & 以 Nexys A7-100T 为目标平台，形成 \texttt{impl50}、报告与 bring-up 路径的统一工程入口 \\
        当前待补齐赛题项 & IF 前端的更完整预取/预测机制、简单应用与演示材料、源码打包与注释率、实体板闭环仍需继续推进 \\
        \bottomrule
    \end{tabularx}
\end{table}

\begin{table}[H]
    \centering
    \small
    \caption{设计范围概览}
    \begin{tabularx}{0.95\textwidth}{L{0.22\textwidth} Y}
        \toprule
        设计内容 & 方案说明 \\
        \midrule
        处理器内核 & 采用 IF/ID/EX/MEM/WB 五级顺序流水，覆盖 hazard、redirect、CSR 与 trap 关键机制 \\
        最小 SoC & 由同步指令存储、数据存储、UART、timer 和 DONE 构成，用于承接软件运行与对外观测 \\
        功能验证 & 给出 smoke、\texttt{rv32 baseline 33/33}、\texttt{rv64 baseline 21/21} 以及扩展 \texttt{full-ui} 验证结果 \\
        性能评估 & 给出 CoreMark short、strict \texttt{>=10s} 与 FPGA-like probe 的结果，用于说明性能稳定性与路径一致性 \\
        FPGA 原型适配 & 给出 Vivado \texttt{impl50} 的资源、时序和原型路径固定配置结果，用于说明原型实现可行性 \\
        后续工作 & 实体板连线验证、现场观测和最终板级 I/O 约束确认需在外部条件具备后继续开展 \\
        \bottomrule
    \end{tabularx}
\end{table}

本文重点说明以 \texttt{RV32I + Zicsr} 为冻结性能口径的处理器及其最小 SoC 实现；
同时给出 \texttt{RV32/RV64} 扩展 \texttt{full-ui} 与 baseline 结果，用于说明该工程已具备双 XLEN 验证能力与更广测试范围下的稳定性。
